% eksperymentalne tłumaczenie na włoski

%

\section{Przyrządy potężniejsze niż cyrkiel i linjka bez podziałki} % lang-pl
\section{Strumenti più potenti del compasso e della riga non graduata} % lang-it

\subsection{Konstrukcja neusis} % lang-pl
\subsection{Costruzione per neusis} % lang-it

Konstrukcja neusis z linijką z podziałką. % lang-pl
Costruzione per neusis mediante una riga graduata. % lang-it
%  νεύειν (neuein) 'incline towards'; plural: νεύσεις, neuseis => https://en.wikipedia.org/wiki/Neusis_construction
% https://en.wikipedia.org/wiki/Straightedge_and_compass_construction#Markable_rulers
Archimedes poda konstrukcję neusis dla siedmiokąta foremnego. % lang-pl
Archimede descriverà una costruzione per neusis del settagono regolare. % lang-it

\subsection{Konstrukcja stożkowa} % lang-pl
\subsection{Costruzioni coniche} % lang-it

Oprócz cyrkla i linijki mamy do dyspozycji narzędzie, które kreśli stożkowe o zadanych ognisku, kierownicy i mimośrodzie. % lang-pl
Taką samą moc ma zestaw złożony z cyrkla, linijki i wykresu paraboli $y = x^2$ dla $x \in [0, 1]$ albo sznurka, który pozwala kreślić elipsy o zadanej półosi i ogniskach. % lang-pl
Starożytni Grecy będą wiedzieć, że podwojenie sześcianu oraz trysekcja dowolnego kąta mają konstrukcje stożkowe. % lang-pl
Kwadratura koła nie ma. % lang-pl
Oltre al compasso e alla riga disponiamo di uno strumento capace di tracciare coniche assegnandone fuoco, direttrice ed eccentricità. % lang-it
La stessa potenza costruttiva si ottiene con un compasso, una riga e il grafico della parabola $y=x^2$ per $x \in [0,1]$, oppure con una cordicella che permetta di tracciare ellissi con semiassi e fuochi assegnati. % lang-it
Gli antichi Greci sapevano che la duplicazione del cubo e la trisezione di un angolo arbitrario ammettono costruzioni coniche. % lang-it
La quadratura del cerchio no. % lang-it
% \todofoot{Solid constructions} % https://en.wikipedia.org/wiki/Straightedge_and_compass_construction#Solid_constructions % https://en.wikipedia.org/wiki/Straightedge_and_compass_construction#Angle_trisection_2

\subsection{Origami} % lang-pl
\subsection{Origami} % lang-it

Patrz sekcja \ref{section_origami}. % lang-pl
Si veda la sezione \ref{section_origami}. % lang-it

%